\section{Introduction}\label{sec:intro}

\IEEEPARstart{A}{} cone crusher sets how much work the grinding mills downstream of it must do, and the concentrators of
Chilean copper mining took 55.7\pct{} of the sector's electricity in 2024~\cite{cochilco}. A model that predicts the
throughput and the power of a crusher from its closed-side setting and its feed size is therefore a planning tool of
direct economic weight, and the models in use are of two kinds. Population-balance models in the tradition of
Whiten~\cite{whiten,napier} describe the classification and the breakage of the particles in the chamber and carry a small
number of fitted scalars; empirical regressions of the outputs on the operating variables describe nothing and fit
everything~\cite{yamashita}. Industrial surveys of a crusher are few and costly, and a practitioner with ten of them must
choose between the two.

CrusherBayes~\cite{cb002}, the calibration method of the ChancaDEM studio~\cite{chancadem}, sits between the kinds. It
calibrates the eight fitted scalars of a Whiten engine of a Metso HP500 cone crusher as a Bayesian posterior on the ten
surveys of Rocha et al.~\cite{rocha}, with a structural discrepancy in the sense of Kennedy and O'Hagan~\cite{ko} that is a
plane in the standardized setting and feed size, integrated out exactly. On the ten surveys, held out one at a time, it
beats the engine's power form refit by least squares by 11.63 points of mean absolute percentage error (paired bootstrap
interval $-18.25$ to $-4.47$) and its 80\pct{} bands cover eight of the ten held-out surveys on throughput and nine on power, but it ties the
line in the setting alone: the paired difference, $+1.54$
points with a bootstrap interval of $-1.11$ to $4.10$, cannot be told from zero. That paper proved that the predictive mean
is a shrunk ridge regression on the discrepancy regressors plus the engine's departure from their plane. It left open the
question the data could not answer: under which conditions, with how many surveys and over which operating range, does the
physics in the method pay for itself against an empirical regression?

This paper answers that question theoretically and in simulation, which ten surveys cannot do because the answer depends on
the plant's true mean function, and no finite set of surveys reveals it. We treat the method as an estimator and ask for its
risk. The answer has an exact form. Given the scalars, the expected squared error of the method at any operating point is a
squared bias in three quantities, the plant's departure from a plane in the regressors, the engine's departure from the same
plane and the distance between the two planes, plus the variance of a ridge fit and the noise. The comparison with least
squares then reduces to one inequality: the engine helps only through the part of it that is not a plane, and only when that
part points the same way as the plant's own curvature and is not more than twice as large. The rest of the paper follows from
this.

The contributions are:
\begin{enumerate}
\item an exact risk decomposition of the calibrated predictor given its scalars (Theorem~\ref{thm:risk}), with the criterion
for beating an empirical regression in the same regressors (Corollary~\ref{cor:align}), the statement that an engine affine
in the regressors is only a shrinkage target (Corollary~\ref{cor:affine}), and the exact coverage of its band
(Corollary~\ref{cor:coverage});
\item a fixed-design decomposition into a curvature gain that persists and a plane gain that decays with the number of
surveys, with the risk-optimal scale of the discrepancy prior (Theorem~\ref{thm:fixed}) and the crossing sample size against
the setting-only line (Proposition~\ref{prop:line});
\item the relation of the Kennedy and O'Hagan prior to the orthogonal prior of Plumlee~\cite{plumlee} and the scaled prior of
Gu and Wang~\cite{guwang} in this finite-dimensional case (Proposition~\ref{prop:priors}), and an information ceiling that no
number of surveys at a fixed speed can break (Proposition~\ref{prop:bound});
\item a simulation study with known truths, designed and reported by the ADEMP structure of Morris et al.~\cite{morris}, on
virtual plants built from the engine and from the published capacity law of Rocha et al.~\cite{rocha}, with every
performance number given with its Monte Carlo standard error and every theoretical risk checked against its simulation;
\item a design study of the information that added outputs, added factors and added surveys bring to the scalars, and a
simulation-based calibration~\cite{talts,modrak} of the sampler that the method relies on.
\end{enumerate}

The main results are stated in the abstract and in Section~\ref{sec:results} with their uncertainty. In brief, the engine
cannot improve on a regression for throughput at all on this machine, because its capacity law is affine in the setting; it
can on a plant whose capacity curves, if the engine's capacity curves the same way, and the size of that gain is set by the
alignment of the two curvatures and not by the number of surveys. The answer to the practitioner's question is therefore a
property of the engine and of the plant's operating range, which the theory makes computable, rather than of the size of the
set of surveys.

Section~\ref{sec:related} places the work among the published studies of calibration, of ridge regression under
misspecification, of design for calibration and of the evaluation of samplers. Section~\ref{sec:problem} states the problem
and the notation; Section~\ref{sec:method} the method and its closed forms; Section~\ref{sec:theory} the theory;
Section~\ref{sec:design} the study design and its hypotheses; Section~\ref{sec:results} the results;
Section~\ref{sec:ablation} the ablation; Sections~\ref{sec:discussion} and~\ref{sec:limitations} the discussion and the
limitations; and Sections~\ref{sec:repro} and~\ref{sec:conclusions} the reproducibility record and the conclusions.
